% Chương 6 -- Kết luận và đề xuất.
\chapter{Kết luận và đề xuất}\label{ch:conclusion}

\section{Kết luận}

Trên tập dữ liệu đơn giao đúng hạn của Olist, Random Forest đạt PR-AUC mức đơn 0,146502 và Recall@5\% 0,121225 qua ba fold phát triển. PR-AUC ở đây đo khả năng đưa các đơn có review 1--2 sao lên đầu hàng đợi; Recall@5\% đo tỷ lệ các đơn đó xuất hiện trong 5\% đơn có score cao nhất. Vì vậy, khi nhân viên chỉ xem 5\% đơn đầu hàng đợi, họ gặp khoảng 12,1\% trường hợp review thấp, thay vì khoảng 5\% nếu chọn ngẫu nhiên. Kết quả cho thấy score có thể hỗ trợ ưu tiên review, nhưng không cho phép tự động quyết định cách xử lý hay kết luận rằng sự hài lòng của khách hàng đã được cải thiện.

Để score này đi vào vận hành mà không bị diễn giải quá mức, hệ thống lưu nó trong snapshot, tức bản ghi cố định tại thời điểm chấm điểm. Khi một case được nhận, liên lạc hay đóng, workflow chỉ cập nhật lịch sử xử lý; score ban đầu không bị sửa. Cùng với phân quyền, SLA và audit, cách tách này cho phép kiểm tra ai đã xử lý case nào, lúc nào và dựa trên thông tin nào. Do đó, mô hình chỉ tạo thứ tự cần xem xét, còn hành động với khách hàng vẫn thuộc về nhân viên.

\section{Ý nghĩa của kết quả}

Kết quả trên chỉ có nghĩa trong các ranh giới đã đặt ra. Thứ nhất, mô hình chỉ dùng dữ liệu biết tại thời điểm dự báo; vì vậy nó không học từ thông tin xuất hiện sau đó. Thứ hai, dữ liệu được chia theo thời gian để mô phỏng tình huống dùng quá khứ dự báo giai đoạn sau. Thứ ba, review là nhãn của đơn hàng, nên việc gán nhãn này cho từng item không cho phép kết luận item nào gây ra bất mãn. Cuối cùng, các metric dùng để chọn mô hình không thể đồng thời được xem là kết quả xác nhận. Những ranh giới này không làm mô hình mạnh hơn, nhưng ngăn việc suy rộng từ một bài toán xếp hạng sang một giải thích nhân quả hoặc một cam kết vận hành.

\section{Hạn chế và hướng phát triển}\label{sec:future-work}

Hạn chế quan trọng nhất nằm ở trạng thái của bằng chứng. Historical test period đã được quan sát trong khi phát triển, nên không còn là một holdout hợp lệ. Vì vậy, các metric hiện tại chỉ là kết quả thăm dò. Báo cáo cũng chưa có khoảng tin cậy bootstrap theo cụm đơn, đường PR và đường hiệu chỉnh trên holdout, dự đoán mức dòng để kiểm tra hoặc giá trị SHAP để giải thích đóng góp của đặc trưng. Bên cạnh đó, dữ liệu Olist chỉ mô tả giai đoạn 2016--2018 và không ghi nhận một can thiệp CSKH có đối chứng. Từ dữ liệu này không thể suy ra hàng đợi đã làm thay đổi hành vi hay đánh giá của khách hàng.

Hướng phát triển nên theo thứ tự sau:
\begin{enumerate}
  \item Trước hết, khóa một replacement holdout chưa được quan sát. Model, calibration, cách tổng hợp điểm và ngưỡng phải được chốt trước khi mở tập này; sau đó mới xuất metric, khoảng tin cậy và artefact dự đoán để chuyển từ kết quả thăm dò sang kết quả xác nhận.
  \item Khi mô hình đã qua bước xác nhận, hoàn thiện đường chấm điểm theo lô bằng model artifact thật, rồi ghi phiên bản snapshot và ngưỡng vào schema \texttt{ml}. Chỉ sau bước đó mới thay dữ liệu demo trong ứng dụng. SHAP hay một cơ chế giải thích khác chỉ nên xuất hiện trên giao diện khi có artefact tương ứng và hướng dẫn tránh diễn giải nhân quả.
  \item Tiếp theo, thu thập dữ liệu vận hành sau dự báo, gồm lần liên lạc, kết quả, thời gian phản hồi và đánh giá sao sau đó. Nếu muốn đo hiệu quả CSKH, cần có nhóm đối chứng hoặc cách phân bổ ngẫu nhiên phù hợp; dữ liệu quan sát thuần túy không đủ để suy tác động.
  \item Cuối cùng, chỉ mở rộng đặc trưng bằng dữ liệu thực sự biết tại thời điểm giao hàng hoặc thiết kế lại bài toán ở cấp đơn. Việc tiếp tục tinh chỉnh thuật toán trên cùng dữ liệu chỉ có ý nghĩa khi đi kèm giả thuyết và giao thức mới.
\end{enumerate}
